\documentclass[10pt,a4paper]{amsart}
\usepackage[margin=19mm]{geometry}
\usepackage{amsmath,amssymb,mathtools}
\usepackage[scr=rsfs,cal=euler]{mathalfa}
\usepackage{xfrac}
\usepackage[unicode,hidelinks,bookmarks=false]{hyperref}
\newcommand{\ie}{i.e.\ }
\newcommand{\cf}{cf.\ }
\newcommand{\resp}{respectively\ }
\newcommand{\Subsection}{\subsection}
\providecommand{\nobreakdash}{\nobreak\hbox{-}}
\setlength{\emergencystretch}{2em}
\multlinegap0pt
\usepackage{bm}
\usepackage{bbm}
\usepackage{dsfont}
\usepackage{xspace}
\usepackage{cancel}
\usepackage{mathtools}
\usepackage{amsthm}
\makeatletter
\@ifundefined{subsection}{\newtheorem{subsection}{Subsection}}{}
\@ifundefined{prop}{\newtheorem{prop}[subsection]{Proposition}}{}
\makeatother
\DeclareMathOperator{\Ig}{Ig}
\DeclareMathOperator{\dvsr}{div}
\DeclareMathOperator{\ord}{ord}
\DeclareMathOperator{\spec}{Spec}
\newcommand{\EE}{\mathcal{E}}
\newcommand{\OO}{\mathcal{O}}
\renewcommand{\and}{\quad\text{and}\quad}
\title[New unlikely intersections on elliptic surfaces]{New unlikely intersections on elliptic surfaces}
\author{Douglas Ulmer, Jos\'e Felipe Voloch}
\date{Source: arXiv:2508.06680v1 (CC BY 4.0)}
\begin{document}
\maketitle

\noindent\emph{Source and licence.} New unlikely intersections on elliptic surfaces. Version arXiv:2508.06680v1 (CC BY 4.0), distributed under the Creative Commons Attribution 4.0 International licence, as recorded in the arXiv licence metadata for this exact version and shown on the abstract page https://arxiv.org/abs/2508.06680v1. Full rights notice: licenses/arxiv-2508.06680-CC-BY-4.0.txt.\par\smallskip

\noindent\textit{Editorial scope.} Counted unit: the Prop whose source label is \texttt{prop:local-points-and-bounds} in the article (source0.tex lines 918--934, source label \texttt{prop:local-points-and-bounds}), delivered together with the article's own complete original proof of it (source0.tex lines 936--1000, one \texttt{proof} environment of 65 lines). No mathematics is written, repaired or re-derived: statement and proof are the article's own text line for line. Required context retained uncounted: eq:Broumas2 (lines 890--895); eq:divs1 (lines 905--907); eq:divs2 (lines 909--911). Every definition, display and label of those lines is retained unchanged. Omitted from this package and not redistributed: 1--889, 896--904, 908--908, 912--917, 935--935, 1001--1514. Nothing inside a retained excerpt is omitted or altered: every retained body is delivered exactly as the article's lines are written, so no omission receipt of any kind accompanies this package. Citations: the retained excerpts carry no citation command at all, so no bibliography entry of the article is transcribed and no reference is redistributed. Reference adaptation: none was needed and none was made; every reference inside the delivered unit resolves inside it, so no unresolved reference or citation is produced. Printed slips kept literal: no printed word, symbol, hypothesis or number of the counted statement or of its proof was altered. Layout and class: the article's own class is \texttt{amsart} with packages including amssymb, amscd, amsmath, amsrefs, xy, fullpage, fontenc, libertine, stackengine, mathtools, graphicx, multirow, colortbl, pdfsync; the wrapper is the lane's pinned \texttt{amsart} editorial wrapper at 10pt on A4, which restates the theorem-like environments the retained text needs and adds the article's own macro definitions that the retained text needs (\texttt{\textbackslash EE}, \texttt{\textbackslash Ig}, \texttt{\textbackslash OO}, \texttt{\textbackslash and}, \texttt{\textbackslash dvsr}, \texttt{\textbackslash ord}, \texttt{\textbackslash spec}), with extra wrapper packages (bm, bbm, dsfont, xspace, cancel, mathtools). The delivered numbering is the unit's own. Assets: no retained span contains a figure environment, a graphics-inclusion command, a PDF-page inclusion, a table or a TikZ picture; no image file is copied into this package, the package declares no asset dependency and no image file is redistributed with it. Licence: version arXiv:2508.06680v1 of this article is distributed under the Creative Commons Attribution 4.0 International licence, as recorded in the arXiv licence metadata for this exact version and shown on the abstract page https://arxiv.org/abs/2508.06680v1. A full rights notice is delivered at licenses/arxiv-2508.06680-CC-BY-4.0.txt. Characters: every retained excerpt is pure ASCII, so the delivered engine, XeLaTeX with its default Latin Modern text font, reports no missing character.
\begin{equation}\label{eq:Broumas2}
\lambda=\alpha^{p-2}\frac{dx'(Q_0)}{2y'(Q_0)}
  \and
  \mu(P)=\alpha^p\wp\left(
    \frac{dx'(Q)/2y'(Q)}{dx'(Q_0)/2y'(Q_0)}\right)  
\end{equation}

\begin{equation}\label{eq:divs1}
\dvsr(\alpha)=(ss)\and\dvsr(\lambda)=(p-2)(ss)-(cusps),  
\end{equation}

\begin{equation}\label{eq:divs2}
  \dvsr\left(dx'(Q_0)/2y'(Q_0)\right)=-(cusps).
\end{equation}

\begin{prop}\label{prop:local-points-and-bounds}\mbox{}
  \begin{enumerate}
  \item   Let $z\in\EE$ be a closed point, and let $t=\pi(z)$.  Then
    there is a section $R$ of $\pi$ over $\spec\OO_t$ such that $pR(t)=z$.
  \item If $R$ is a local section of $\pi$ over $\spec\OO_t$ with
    $pR(t)=P(t)$ and if $\iota=(pR,P)_t$ is the local intersection
    number, then we have
\[\ord_t(\nu(P))\ge
  \begin{cases}
    \iota-1&\text{if $t$ is not a supersingular point,}\\
    p(\iota-1)-(p-1)(p-2)
       &\text{if $t$ is supersingular and $\iota< p$,}\\
    \iota+p-2
       &\text{if $t$ is supersingular and $\iota\ge p$.}
  \end{cases}\]
  \end{enumerate}
\end{prop}

\begin{proof}
  (1) Because $\EE$ and $\EE'$ have everywhere semistable reduction,
  the connected components of the identity
  $G=\left(\pi^{-1}(t)\right)^0$ and are either elliptic curves or
  multiplicative groups, and multiplication map $p:G\to G$ is
  surjective.  Moreover, since $\Ig_1(N)\to X_1(N)$ is completely
  split over the cusps, the component groups of $\EE$ have order
  dividing $N$ and thus of order prime to $p$.  Therefore, $p$ is also
  surjective on the component groups, and using the snake lemma, on
  all of $\EE$.  Suppose that $w\in \EE$ satisfies $pw=z$.  Since
  $\pi$ is smooth, by Hensel's lemma $w$ lifts to a point $R$ with
  values in $\OO_{t}$, and we have $pR(t)=pw=z=P(t)$, as desired. 

  (2) With $P$ and $R$ as in the statement, let $R_0=P-pR$, a local
  section of $\pi$ which specializes to the identity at $t$.  Since
  $\nu$ kills $pE(K_t)$, we have $\nu(P)=\nu(R_0)$.  Moreover,
  \[\iota:=(pR,P)_t=(O,R_0)_t,\] %
  so we are reduced to showing that if $R_0$ is a local section of
  $\pi$ which specializes to the identity and if $\iota=(O,R_0)_t$,
  then $\nu(R_0)$ is bounded below by the quantity displayed in the
  statement of part (2).

  If $t$ is not a supersingular point, then $V$ is \'etale over $t$
  and we may choose $Q_0\in E^{(p)}(\OO_t)$ with $VQ_0=R_0$ and with
  $Q_0$ specializing to the identity at $t$.  With this choice, $Q_0$
  must meet the identity section to order $\iota$, i.e., we have
  $\ord_t(x'(Q_0))=-2\iota$ and $\ord_t(y'(Q_0))=-3\iota$.  We find
  that
\[\ord_t(dx'(Q)/2y'(Q))\ge \iota-1,\]
and formulas \eqref{eq:Broumas2}, \eqref{eq:divs1}, and
\eqref{eq:divs2} then imply that
\[\ord_t(\nu(P))\ge \iota-1.\]

Now assume that $t$ is supersingular.  Then the points $Q_0$ with
$VQ_0=R_0$ are defined over an extension $\OO'$ of $\OO_t$ which may
be ramified, and in order to compute the different, we consider
$V:E^{(p)}\to E$ in the formal groups of $E$ and $E^{(p)}$.  It takes
the shape
\[ V(z)=Az+uz^p+\cdots\] %
where $u$ is a unit at $t$.  Since $R_0$ meets the zero section to
order $\iota$, the choice of $Q_0$ meeting the zero section
corresponds to a solution of
\[v\pi_t^\iota=Az+uz^p+\cdots,\] %
where $v$ is a unit.

Since $A$ vanishes to order $p-1$, a consideration of slopes shows
that if $\iota$ is $<p$, $\OO'$ is (wildly) ramified over $\OO$ with
ramification break $m=p-\iota$.  We find that $dx'(Q_0)/2y'(Q_0)$
vanishes in $\Omega^1_{\OO'/k}$ to order $(p-1)(p-\iota+1)$, and
$Q_0$ meets the zero section to order $\iota$.  Using
\eqref{eq:Broumas2}, \eqref{eq:divs1}, and \eqref{eq:divs2}, we find
that
\begin{align*}
  \ord_t(\nu(P))&\ge (2p-2)+(\iota-1)-(p-1)(p-\iota+1)\\
  &=p(\iota-1)-(p-1)(p-2).
\end{align*}

If $\iota\ge p$, then the choice of $Q_0$ meeting the zero section is
defined over $\OO_t$ and meets the zero section to order
$\iota-(p-1)$, and using \eqref{eq:Broumas2}, \eqref{eq:divs1}, and
\eqref{eq:divs2}, we find that
\[\ord_t(\nu(P))\ge \iota+p-2.\]

This completes the proof of the proposition.
\end{proof}

\end{document}
